\chapter{Quantum Channels and Positive Maps: Supporting Results}
\label{ch:channels_supporting_results}

This appendix supports Chapter~\ref{ch:channels}. It collects the
supporting results for complete positivity, density matrices, mean-ergodic
projections, weighted traces, and idempotent retractions.

\section{Kraus representations and complete positivity}

This section proves the two complete-positivity consequences stated after
Definition~\ref{def:cp_map}: entrywise positivity on block matrices and the
associated completely positive map between matrix $C^*$-algebras.

\begin{theorem}[Entrywise complete positivity from a Kraus representation]
    \label{thm:cp_map_entrywise_complete_positive}
    \lean{IsCPMap.map_cstarMatrix_nonneg}
    \leanok
    \uses{def:cp_map}
    Let $E : \MN{D} \to \MN{D}$ be completely positive in the Kraus sense.
    For every $k$ and every positive block matrix
    $M \in M_k(\MN{D})$, the entrywise image
    $(E(M_{ab}))_{a,b}$ is again positive in $M_k(\MN{D})$.
\end{theorem}

\begin{proof}\leanok
    Write $E$ in the Kraus form (\ref{eq:channel_kraus_representation}).
    For each $i$, let $d_i$ be the block-diagonal matrix with $K_i$ on every
    diagonal block. Then
    \begin{align}
        (E(M_{ab}))_{a,b}
         & =\sum_i d_i M d_i^\dagger.
        \notag
    \end{align}
    Each term on the right is positive because conjugation preserves
    positivity, and a finite sum of positive matrices is positive.
\end{proof}

\begin{theorem}[Kraus maps as abstract completely positive maps]
    \label{thm:cp_map_to_abstract_completely_positive}
    \lean{IsCPMap.toCompletelyPositiveMap}
    \leanok
    \uses{def:cp_map, thm:cp_map_entrywise_complete_positive}
    Every Kraus-represented completely positive map
    $E : \MN{D} \to \MN{D}$ determines a completely positive map between
    the matrix $C^*$-algebras in the entrywise positivity sense.
\end{theorem}

\begin{proof}\leanok
    The defining entrywise positivity condition is exactly
    Theorem~\ref{thm:cp_map_entrywise_complete_positive}.
\end{proof}

Thus the Kraus formulation of Definition~\ref{def:cp_map} has the two
consequences cited there: entrywise complete positivity and the associated
abstract completely positive map.

\section{Density matrices and the Ces\`aro construction}

This section proves the compactness, convexity, nonemptiness, and invariance
facts used in Theorem~\ref{thm:cesaro_fixedpoint}. Write $\mc{D}_D$ for the
density matrices of Definition~\ref{def:density_matrices}; the Ces\`aro means
$\sigma_N$ and their telescope identity remain in
Definition~\ref{def:cesaro_mean} and Theorem~\ref{thm:cesaro_telescope}.

\begin{theorem}[Density matrices are compact]\label{thm:density_compact}
    \lean{densityMatrices_isCompact}
    \leanok
    \uses{def:density_matrices}
    $\mc{D}_D$ is compact in the entrywise topology on $\MN{D}$.
\end{theorem}

\begin{proof}
    \leanok
    The positive semidefinite cone is closed, being the preimage of the closed
    non-negative cone under continuous quadratic forms. If $\rho \ge 0$ with
    $\tr(\rho) = 1$, then each entry satisfies
    $|\rho_{ij}|^2 \le \rho_{ii}\rho_{jj} \le 1$: the diagonal entries are
    non-negative and sum to~$1$, and the off-diagonal bound is the positive
    semidefinite Cauchy--Schwarz inequality. Hence the matrix entries are
    uniformly bounded, and Heine--Borel gives compactness.
\end{proof}

\begin{theorem}[Density matrices are convex]\label{thm:density_convex}
    \lean{densityMatrices_isConvex}
    \leanok
    \uses{def:density_matrices}
    $\mc{D}_D$ is convex.
\end{theorem}

\begin{proof}\leanok
    A convex combination of positive semidefinite matrices is positive
    semidefinite, and the trace is linear.
\end{proof}

\begin{theorem}[Density matrices are nonempty]\label{thm:density_nonempty}
    \lean{densityMatrices_nonempty}
    \leanok
    \uses{def:density_matrices}
    For $D \ge 1$, $\mc{D}_D \neq \varnothing$.
\end{theorem}

\begin{proof}\leanok
    $\frac{1}{D}\Id_D$ is a density matrix.
\end{proof}

\begin{theorem}[Channels preserve density matrices]\label{thm:channel_preserves_density}
    \lean{IsChannel.map_densityMatrices}
    \leanok
    \uses{def:channel, def:density_matrices}
    If $E$ is a quantum channel, then
    $E(\mc{D}_D) \subseteq \mc{D}_D$.
\end{theorem}

\begin{proof}\leanok\uses{thm:cp_is_positive}
    Complete positivity implies positivity by
    Theorem~\ref{thm:cp_is_positive}, so $E(\rho) \ge 0$; trace preservation
    gives $\tr(E(\rho)) = \tr(\rho) = 1$.
\end{proof}

Together these results keep the Ces\`aro orbit inside a compact nonempty set of
density matrices, as required in Theorem~\ref{thm:cesaro_fixedpoint}.

\section{Preservation under the mean-ergodic projection}

This section proves the positivity, trace-preservation, and unitality assertions
used in Theorem~\ref{thm:positive_trace_preserving_mean_ergodic_projection}.
For a bounded-orbit endomorphism $T$, let $P_T$ denote the mean-ergodic
projection of Definition~\ref{def:mean_ergodic_projection}.

\begin{theorem}[Positive mean-ergodic projection]
    \label{thm:positive_mean_ergodic_projection}
    \lean{IsPositiveMap.meanErgodicProjection_isPositiveMap}
    \leanok
    \uses{def:positive_map, def:mean_ergodic_projection,
        thm:mean_ergodic_bounded_orbits}
    Let $T:\MN{D}\to\MN{D}$ be positive and have bounded orbits. Then its
    mean-ergodic projection $P_T$ is positive.
\end{theorem}

\begin{proof}\leanok
    Every Ces\`aro average is positive. Therefore, for $X\ge0$, closedness of
    the positive cone and convergence in
    (\ref{eq:channel_mean_ergodic_limit}) give $P_T(X)\ge0$.
\end{proof}

\begin{theorem}[Trace-preserving mean-ergodic projection]
    \label{thm:trace_preserving_mean_ergodic_projection}
    \lean{IsTracePreservingMap.meanErgodicProjection_isTracePreservingMap}
    \leanok
    \uses{def:trace_preserving, def:mean_ergodic_projection,
        thm:mean_ergodic_bounded_orbits}
    Let $T:\MN{D}\to\MN{D}$ be trace-preserving and have bounded orbits.
    Then its mean-ergodic projection $P_T$ is trace-preserving.
\end{theorem}

\begin{proof}\leanok
    Every nonempty Ces\`aro average preserves the trace. Convergence of the
    averages and continuity of the trace therefore give
    $\tr(P_T(X))=\tr(X)$.
\end{proof}

\begin{theorem}[Identity under the mean-ergodic projection]
    \label{thm:unital_mean_ergodic_projection}
    \lean{LinearMap.HasBoundedOrbits.meanErgodicProjection_one}
    \leanok
    \uses{def:mean_ergodic_projection, thm:mean_ergodic_bounded_orbits}
    Let $T:\MN{D}\to\MN{D}$ have bounded orbits. If $T(\Id)=\Id$, then
    $P_T(\Id)=\Id$.
\end{theorem}

\begin{proof}\leanok
    The identity is a fixed point of $T$, so the fixed-point characterization
    of the mean-ergodic projection gives the conclusion.
\end{proof}

These three preservation statements give the corresponding properties of the
Ces\`aro projection in
Theorem~\ref{thm:positive_trace_preserving_mean_ergodic_projection}.

\section{Trace adjoints of ergodic projections}

This section proves the two adjoint facts used in
Theorem~\ref{thm:positive_unital_adjoint_mean_ergodic_retraction}. For a linear
map $S$, write $S^*$ for its adjoint with respect to the trace pairing; for a
bounded-orbit map $T$, write $P_T$ for its mean-ergodic projection.

\begin{theorem}[Trace preservation and the trace adjoint]
    \label{thm:trace_preserving_iff_trace_adjoint_unital}
    \lean{isTracePreservingMap_iff_traceAdjointMap_one}
    \leanok
    \uses{def:trace_preserving, def:trace_pairing_adjoint}
    On any finite full matrix algebra, a linear endomorphism $S$ is
    trace-preserving if and only if its trace-pairing adjoint fixes the
    identity:
    \begin{align}
        (\forall X,\ \tr(S(X))=\tr(X))
         & \quad\Longleftrightarrow\quad S^*(\Id)=\Id.
        \label{eq:channel_tp_adjoint_unital}
    \end{align}
\end{theorem}

\begin{proof}\leanok
    \uses{thm:trace_pairing_adjoint_identity, thm:trace_nondeg}
    The trace-pairing identity gives
    $\tr(S^*(\Id)X)=\tr(S(X))$. The equivalence
    (\ref{eq:channel_tp_adjoint_unital}) follows from nondegeneracy of the
    trace pairing.
\end{proof}

\begin{theorem}[Adjoint mean-ergodic projection]
    \label{thm:adjoint_mean_ergodic_projection}
    \lean{LinearMap.HasBoundedOrbits.traceAdjointMap_meanErgodicProjection_apply_eq_self_iff,
        LinearMap.HasBoundedOrbits.traceAdjointMap_meanErgodicProjection_apply_self,
        LinearMap.HasBoundedOrbits.range_traceAdjointMap_meanErgodicProjection}
    \leanok
    \uses{def:mean_ergodic_projection, thm:mean_ergodic_bounded_orbits}
    Let $T:\MN{D}\to\MN{D}$ have bounded orbits, let $P_T$ be its
    mean-ergodic projection, and let $P_T^*$ be the trace-pairing adjoint of
    $P_T$. Then
    \begin{align}
        (P_T^*)^2     & =P_T^*, \notag  \\
        \range(P_T^*) & =\ker(T^*-\Id).
        \label{eq:channel_adjoint_mean_ergodic}
    \end{align}
    and $P_T^*(Y)=Y$ if and only if $T^*(Y)=Y$.
\end{theorem}

\begin{proof}\leanok
    \uses{thm:trace_pairing_adjoint_identity, thm:trace_nondeg,
        thm:mean_ergodic_bounded_orbits}
    Since $P_T^2=P_T$, the trace-pairing identity gives
    $(P_T^*)^2=P_T^*$. If $T^*(Y)=Y$, decompose
    $X-P_T(X)=(T-\Id)(Z)$. Then
    \begin{align}
        \tr(Y(X-P_T(X)))
         & =\tr((T^*(Y)-Y)Z)=0.
        \notag
    \end{align}
    Nondegeneracy of the trace pairing gives $P_T^*(Y)=Y$. Conversely,
    $P_TT=P_T$ implies $T^*P_T^*=P_T^*$, so every fixed point of $P_T^*$ is a
    fixed point of $T^*$:
    \begin{align}
        P_T^*(Y)=Y
         & \quad\Longrightarrow\quad
        T^*(Y)=T^*(P_T^*(Y))=P_T^*(Y)=Y.
        \notag
    \end{align}
    Hence
    \begin{align}
        Y\in\range(P_T^*)
         & \quad\Longleftrightarrow\quad P_T^*(Y)=Y
        \quad\Longleftrightarrow\quad T^*(Y)=Y
        \quad\Longleftrightarrow\quad Y\in\ker(T^*-\Id).
        \notag
    \end{align}
    The first equivalence uses idempotence, and the second is the fixed-point
    equivalence proved above. This proves
    (\ref{eq:channel_adjoint_mean_ergodic}).
\end{proof}

The first theorem gives unitality of $P_T^*$ from trace preservation of $P_T$;
the second identifies its range and fixed points, completing the adjoint
retraction argument in
Theorem~\ref{thm:positive_unital_adjoint_mean_ergodic_retraction}.

\section{Weighted traces and idempotent retractions}

This section proves the weighted-trace lemma used in
Theorem~\ref{thm:fixedPointsStarSubalgebra} and the three absorption identities
complementing Theorem~\ref{thm:choi_effros_sandwich}. In the absorption
statements, $E$ is a unital idempotent Kraus map whose adjoint has a positive
definite fixed point, as in
Theorem~\ref{thm:map_mem_multiplicativeDomain_of_idempotent}.

\begin{lemma}[Faithfulness of a positive-definite weighted trace]
    \label{lem:posDef_weighted_trace_faithful}
    \lean{Matrix.posSemidef_eq_zero_of_posDef_trace_mul_eq_zero}\leanok
    If $M\succeq0$, $\rho>0$, and $\tr(\rho M)=0$, then $M=0$.
\end{lemma}

\begin{proof}\leanok
    Write $\rho=S^\dagger S$ with $S$ invertible. The matrix
    $SMS^\dagger$ is positive semidefinite and
    $\tr(SMS^\dagger)=\tr(S^\dagger SM)=\tr(\rho M)=0$.
    A positive semidefinite matrix of trace zero vanishes, so
    $SMS^\dagger=0$. Since $S$ and $S^\dagger$ are invertible,
    $SMS^\dagger=0$ implies $MS^\dagger=0$, and hence $M=0$.
\end{proof}

\begin{theorem}[Choi--Effros left absorption]
    \label{thm:choi_effros_left}
    \lean{Kraus.choi_effros_left}\leanok
    \uses{def:kraus_fixed_points, def:kraus_map}
    Under the same hypotheses, for all $X,Y\in\MN{D}$,
    $E(E(X)Y)=E(E(X)E(Y))$.
\end{theorem}

\begin{proof}\leanok
    \uses{thm:map_mem_multiplicativeDomain_of_idempotent}
    Theorem~\ref{thm:map_mem_multiplicativeDomain_of_idempotent} places
    $E(X)$ in the multiplicative domain, so
    $E(E(X)Y)=E(E(X))E(Y)=E(X)E(Y)$.
    Applying the same multiplicativity to $E(X)$ and $E(Y)$ gives
    $E(E(X)E(Y))=E(E(X))E(E(Y))=E(X)E(Y)$ by idempotence.
\end{proof}

\begin{theorem}[Choi--Effros right absorption]
    \label{thm:choi_effros_right}
    \lean{Kraus.choi_effros_right}\leanok
    \uses{def:kraus_fixed_points, def:kraus_map}
    Under the same hypotheses, for all $X,Y\in\MN{D}$,
    $E(XE(Y))=E(E(X)E(Y))$.
\end{theorem}

\begin{proof}\leanok
    \uses{thm:map_mem_multiplicativeDomain_of_idempotent}
    Theorem~\ref{thm:map_mem_multiplicativeDomain_of_idempotent} places
    $E(Y)$ in the multiplicative domain, so
    $E(XE(Y))=E(X)E(E(Y))=E(X)E(Y)$.
    Applying the same multiplicativity to $E(X)$ and $E(Y)$ gives
    $E(E(X)E(Y))=E(E(X))E(E(Y))=E(X)E(Y)$ by idempotence.
\end{proof}

\begin{theorem}[Choi--Effros symmetric absorption]
    \label{thm:choi_effros_left_eq_right}
    \lean{Kraus.choi_effros_left_eq_right}\leanok
    \uses{def:kraus_fixed_points, def:kraus_map}
    Under the same hypotheses, for all $X,Y\in\MN{D}$,
    $E(E(X)Y)=E(XE(Y))$.
\end{theorem}

\begin{proof}\leanok
    \uses{thm:choi_effros_left, thm:choi_effros_right}
    By Theorems~\ref{thm:choi_effros_left}
    and~\ref{thm:choi_effros_right}, both sides equal
    $E(E(X)E(Y))$.
\end{proof}

The left and right identities identify both absorbed products with the
projected product, and the symmetric identity equates them.
